\section{Limitations}

\paragraph{Environment and scope.} The network is synthetic: a controlled
instrument built to isolate mechanisms, not a model of any city. Corridor demand
is a single origin--destination pair choosing among three corridors; the four
minor streets couple the routing decision to traffic that is not routed, but
many simultaneously-routed pairs, with the interactions between their choices,
are outside what is tested. No claim of real-world validation is made anywhere
in this paper. What is portable is the method rather than the numbers: the
descriptors are dimensionless, so a boundary is stated as a degree of saturation
and a controlled-flow share, both measurable on a real network.

\paragraph{Decision context and compliance.} The selector is given the operating
conditions; a deployed system would have to forecast them and would inherit the
forecast's error on top of the selector's. Guided vehicles comply fully with the
policy, and unguided vehicles never respond to conditions at all; real behaviour
lies between these. The reported decision quality is therefore an upper bound on
what the same selector would achieve behind a forecast and partial compliance.

\paragraph{Resolution and statistical reach.} The experiment samples a finite
grid with \NumSeeds{} seeds. Boundaries are reported as brackets between
adjacent sampled levels and never as continuous thresholds, and no refined
switching step is seed-resolved (Section~5.9); resolving the campaign-wide
adaptation effect would need roughly \NumResolvSeedsNeeded{} seeds per context.
Where a factor is associated with a change of preferred policy, that is an
association within a designed factorial, not an identified causal mechanism.

\paragraph{Protocol.} Two deviations are on the record and are described in
Section~\ref{sec:integrity}: the screening gate G1 failed and the campaign proceeded under
deviation D-1, and one declared shift split was infeasible and was replaced by
the nearest feasible construction. Neither is presented as neutral.

\paragraph{Uncertainty machinery.} Split conformal coverage requires
exchangeability, which covariate shift breaks; the conformal gate is calibrated
in-distribution and heuristic outside it, the support gate detects only
extrapolation visible in the feature space, and neither detects concept shift.
Split O8 is a demonstrated instance of a shift the support gate does not see.

\paragraph{Frozen parameters and the emission model.} The policy parameters were
fixed before the experiment and are not optima; a better-tuned load-balancing
tolerance improves that policy by \NumEpsGainPct{}, which strengthens the
central conclusion but means the reported value is not the best available.
\CO{} is computed by the simulator's HBEFA3 implementation for one vehicle
class, without fleet composition, cold starts or gradient, so C2 should be read
as a consistent comparative index rather than an absolute emissions estimate.
Finally, no monetary valuation is used anywhere, which keeps the study free of
invented constants but also means it cannot express a trade-off in currency.
